% Source: https://artofproblemsolving.com/community/c3h2500198_infinitely_nested_jichael_mordans
% Stable identity: aops-topic-2500198-q01
\textbf{The value of the radical.} Let $a_1=\sqrt x$ and $a_{k+1}=\sqrt{x+a_k}$; the infinite radical means $\lim a_k$. Let $L>0$ be the positive root of $L^2=x+L$. By induction $a_k<L$ (since $a_1=\sqrt x<L$ and $a_{k+1}<\sqrt{x+L}=L$) and $a_k$ is increasing (since $a_2>a_1$, and $a_k>a_{k-1}$ implies $a_{k+1}>a_k$). So $a_k$ converges to a limit $\ell>0$ with $\ell=\sqrt{x+\ell}$, i.e. $\ell^2=x+\ell$, so $\ell=L$.

\textbf{Finding $x$.} With $y=2021^{2021}$ we get $y^2=x+y$, so $x=y(y-1)$.

\textbf{Last two digits.} Modulo $100$, $2021\equiv21$, and by the binomial theorem $21^k=(1+20)^k\equiv1+20k\pmod{100}$, because all later terms contain $20^2=400$. Hence
\[y\equiv1+20\cdot2021=40421\equiv21\pmod{100},\qquad x=y(y-1)\equiv21\cdot20=420\equiv20\pmod{100}.\]
The last two digits of $x$ are $\boxed{20}$.
